\section{Conclusion}

The Carry Protocol represents a pivotal leap forward in the confluence of gaming and advertising within the burgeoning web3 domain. By introducing a structured and transparent framework for integrating ads into the gaming experience, Carry not only unlocks a new revenue stream for developers but also paves the way for advertisers to reach a dynamic and engaged audience.

Central to this is the concept of slots — dynamic spaces within games that can host various types of advertisements. The Carry Protocol's meticulous design ensures that these slots integrate seamlessly with the gaming interface, ensuring that advertising is both organic and engaging.

The robust governance structure proposed ensures the protocol's continuous evolution, adapting to the ever-evolving landscape of gaming and digital advertising. Moreover, the emphasis on data-driven insights, facilitated through detailed metrics and experimentation capabilities, ensures that every stakeholder, from developers and advertisers to gamers, has the tools they need to extract maximum value from the system.

In a broader spectrum of gaming environments, particularly within fully on-chain games, slots predominantly belong to the gaming community or individual players. Consequently, advertising revenue is more intricately allocated to players and developers. There is a compelling rationale to anticipate that this direction represents the future of advertising.

In essence, the Carry Protocol is not just a technological innovation but a holistic ecosystem, fostering collaboration, innovation, and sustainable growth in the web3 gaming advertisement realm. As the digital realm continues to expand and the lines between gaming and real-world experiences blur further, the Carry Protocol is poised to be at the forefront, shaping the future of in-game advertising.
